\section{Related Work}
\label{sec:related}

\subsection{Decision-relevant elicitation}
\citet{regan2009regret} develop regret-based reward elicitation for MDPs,
selecting queries to reduce decision-relevant uncertainty about near-optimal
policies rather than to identify a unique reward; only distinctions that change
the optimal (or near-optimal) policy are worth eliciting.
Their setting uses bound-style queries over reward structure, and related work
extends approximate regret-based elicitation in MDPs~\citep{alizadeh2015approximate}.
Model-Free Preference Elicitation~\citep{modelfree2024} likewise moves the
elicitation target off raw parameters toward reward equivalence classes and
recommendation-quality expected value of information, but in recommender
settings without an MDP policy.
The decision-relevant / regret-based \emph{goal} itself is not novel for this
paper.
We instantiate that goal for preference-based robotics: human trajectory-segment
comparisons under Bradley--Terry likelihoods rather than bound queries; a
learned population prior over previous users rather than searching unrestricted
reward space; and a robotics setting with trajectories.

\subsection{Active preference-based reward learning}
Active preference-based reward learning chooses which trajectory comparison to
ask next.
Classical criteria shrink uncertainty over reward parameters: volume removal over
reward weights~\citep{sadigh2017active}, batch volume
removal~\citep{biyik2018batch}, and information gain that favors answerable
questions~\citep{biyik2019asking}.
Later methods use ensemble disagreement~\citep{hejna2023fewshot}, continuous
query synthesis~\citep{zhou2026preferenceekf}, and effort-aware selective
querying~\citep{muraleedharan2025sparq}.
Closest within PbRL, \citet{ellis2024equivalence} acquire information about
reward \emph{equivalence classes} for a \emph{single user}; our target is
instead the regret of the \emph{executed policy} under a \emph{population prior}.
Outside robotics, EPIG~\citep{epig2023} is the general active-learning version
of the claim that parameter information can be the wrong target, measuring
expected \emph{predictive} mutual information between the query label and future
predictions---not policy regret or Regan-style value of information.
We keep the segment-comparison interface of active PbRL, but score queries by
expected reduction in policy regret under a population prior rather than by
parameter volume, information gain, single-user equivalence acquisition, or
predictive mutual information.

\subsection{Population and pluralistic preference models}
A parallel line models preference diversity as structure rather than noise.
As shown by MaxMin-RLHF (Theorem~1)~\citep{chakraborty2024maxmin}, a single
average reward can fail under diverse preferences---the ``average user'' is
nobody.
Distributional preference learning accounts for hidden
context~\citep{siththaranjan2024dpl}; pluralistic and prototype-based models
personalize reward predictors~\citep{chen2025pal}; and clustering methods learn
representative rewards from diverse preferences~\citep{kim2026prec}.
Closest architecturally, variational preference learning encodes annotations
into a user latent with a latent-conditioned reward~\citep{poddar2024vpl}, but
uses fixed survey queries and leaves active selection as future work.
These methods primarily use population structure to \emph{represent} or align to
users.
We use the population to \emph{choose queries}, not only to model users.

\subsection{Active personalization elsewhere}
Outside robot policy learning, comparison-based active preference learning
personalizes multi-dimensional models in language settings~\citep{oh2025ample},
and active utility-based pairwise sampling improves personalized
recommendations~\citep{boroomand2025utility}.
Neither setting has a downstream policy or a regret target.
Across these neighbors, the cell we claim is an instantiation---policy-regret
scoring of human trajectory comparisons under a population prior---not a new
elicitation paradigm.
